% Style follows ../class2/resume2PT26.tex.
% Build in this directory (run twice for page references):
% pdflatex -interaction=nonstopmode -halt-on-error resume5PT26.tex
% Topics follow the syllabus in ../class1/class1PT26.tex.
% Content adapted from ../lectureNote.tex and ../lectureNotes.tex.
\documentclass[a4paper,10pt,pdflatex,ja=standard,jafont=ipaex-type1,
  magstyle=real,scale=1]{bxjsarticle}
\usepackage{lmodern}
\usepackage{amsmath,amssymb,mathtools}
\usepackage{geometry}
\geometry{reset,left=22mm,right=22mm,top=20mm,bottom=17mm,
  headheight=13pt,headsep=6mm,footskip=10mm}
\usepackage{xcolor}
\usepackage{enumitem}
\usepackage{array,booktabs,tabularx}
\usepackage{tikz}
\usepackage{fancyhdr}
\usepackage{hyperref}
\definecolor{ink}{HTML}{193A5C}
\definecolor{accent}{HTML}{237D82}
\definecolor{soft}{HTML}{EAF3F3}
\definecolor{muted}{HTML}{52616B}
\newcommand{\coursename}{確率論と統計学 I}
\hypersetup{unicode=true,colorlinks=true,urlcolor=accent,linkcolor=ink,
  pdfauthor={中島秀太},pdftitle={第5回：代表的な確率分布}}
\pagestyle{fancy}
\fancyhf{}
\fancyhead[L]{\small\color{ink}\coursename{}・授業要約と演習}
\fancyhead[R]{\small\color{ink}第5回}
\fancyfoot[C]{\small\color{ink}\thepage/\pageref*{lastpage}}
\renewcommand{\headrulewidth}{0.3pt}
\setlength{\parindent}{0pt}
\setlength{\parskip}{2pt}
\setlength{\abovedisplayskip}{4pt}
\setlength{\belowdisplayskip}{4pt}
\setlength{\abovedisplayshortskip}{3pt}
\setlength{\belowdisplayshortskip}{4pt}
\setlist[itemize]{leftmargin=1.8em,itemsep=1pt,topsep=2pt,parsep=0pt}
\setlist[enumerate]{leftmargin=2em,itemsep=3pt,topsep=3pt,parsep=0pt}
\newcommand{\topic}[1]{\par\vspace{5pt}
  {\large\sffamily\mdseries\color{accent}#1}\par\vspace{2pt}}
\newcommand{\handouttitle}[1]{%
  \begin{center}
    {\Large\color{ink}#1}\par\vspace{3mm}
    {\color{accent}\rule{0.88\linewidth}{0.7pt}}
  \end{center}\vspace{-2mm}}
\newcommand{\keybox}[1]{\par\vspace{3pt}\begingroup
  \setlength{\fboxsep}{5pt}
  \colorbox{soft}{\parbox{\dimexpr\linewidth-2\fboxsep\relax}{#1}}%
  \endgroup\par\vspace{2pt}}
\newcommand{\examplelabel}[1]{{\sffamily\mdseries\color{ink}#1}}
\newcommand{\columnheading}[1]{\par\vspace{5pt}
  {\fontsize{11pt}{15pt}\selectfont\sffamily\mdseries\color{accent}#1}\par\vspace{2pt}}
\newcommand{\columnstart}[1]{\clearpage
  \setlength{\parskip}{2.5pt}
  \fancyhead[L]{\small\color{ink}\coursename{}・発展的内容}
  {\small\sffamily\mdseries\color{accent}発展的内容}\par\vspace{4pt}
  {\Large\color{ink}#1}\par\vspace{2pt}
  {\color{accent}\rule{\linewidth}{0.7pt}}\par\vspace{4pt}}
\newcommand{\F}{\mathcal F}
\newcommand{\R}{\mathbb R}
\newcommand{\N}{\mathbb N}
\newcommand{\E}{\mathbb E}
\newcommand{\Var}{\operatorname{Var}}
\newcommand{\Cov}{\operatorname{Cov}}
\newcommand{\ind}{\mathrm{1}}
\newcommand{\iid}{\text{独立同分布}}
\begin{document}
\fontsize{10pt}{13.5pt}\selectfont
\begingroup
\fancyhead[L]{\small\color{ink}\coursename{}・前回の演習解答}
\handouttitle{第4回\quad 演習問題の解答}

\topic{演習1：サイコロの事象の独立性}
公平なサイコロの出目を $D$、$A=\{D\text{ は偶数}\}$、
$B=\{D\le3\}$、$C=\{D\le4\}$ とする。
$$
  P(A)=\frac12,\quad P(B)=\frac12,\quad P(C)=\frac23,
  \qquad P(A\cap B)=\frac16\ne\frac14=P(A)P(B).
$$
よって $A,B$ は独立ではない。一方、$A\cap C=\{2,4\}$ なので
$P(A\cap C)=1/3=P(A)P(C)$ となり、$A,C$ は独立である。

\topic{演習2：無相関でも独立でない和と差}
独立な公平サイコロの出目 $U,V$ に対し、$S=U+V$、$T=U-V$ とする。
$\Var(U)=\Var(V)=35/12$ と共分散の双線形性から
$$
  \Cov(S,T)=\Var(U)-\Var(V)=0.
$$
しかし、$\{S=2,T=0\}=\{U=V=1\}$、$\{T=0\}=\{U=V\}$ より
$$
  P(S=2,T=0)=\frac1{36},\qquad
  P(S=2)P(T=0)=\frac1{36}\cdot\frac16=\frac1{216}.
$$
積公式が破れるので、$S,T$ は独立ではない。

\topic{演習3：有限個の和の分散}
二次モーメントが有限な $X_1,\ldots,X_n$ について、$W_j=X_j-\E[X_j]$ とおく。
$\E|W_iW_j|<\infty$ は Cauchy--Schwarz の不等式から従う。したがって
$$
  \begin{aligned}
    \Var\left(\sum_{j=1}^nX_j\right)
      &=\E\left[\left(\sum_{j=1}^nW_j\right)^2\right]\\
      &=\sum_{j=1}^n\E[W_j^2]+2\sum_{i<j}\E[W_iW_j]
       =\sum_{j=1}^n\Var(X_j)+2\sum_{i<j}\Cov(X_i,X_j).
  \end{aligned}
$$
対独立で各共分散は $0$ になるので十分である。相互独立でも成り立つ。

\topic{演習4：三つの符号の関係}
独立な公平な符号 $X,Y\in\{-1,1\}$ と $Z=XY$ を考える。
$(X,Y)$、$(X,Z)$、$(Y,Z)$ のいずれも、四通りの符号の組が各確率 $1/4$ で現れる。
各変数の周辺確率は $1/2$ なので、どの二つ組も独立である。
一方、$Z=XY$、$X=YZ$、$Y=XZ$ より、どの二つの値からも残りが決まる。
$P(X=Y=Z=1)=1/4\ne1/8$ だから、三つは相互独立ではない。
\clearpage
\endgroup

\handouttitle{第5回\quad 代表的な確率分布}
\keybox{分布を表す式、パラメータの意味、期待値と分散を組にして覚えよう。
回数を数える離散分布と、待ち時間を表す連続分布の関係が見えてくる。}

\topic{確率質量と確率密度}
$X$ の分布を $\mu_X(B)=P(X\in B)$ とする。離散型なら
$$P(X=x)=p_X(x),\qquad p_X(x)\ge0,\qquad \sum_xp_X(x)=1.$$
ある非負可測関数 $f_X$ により $\mu_X(B)=\int_B f_X(x)\,dx$ と書けるとき、
$f_X$ を確率密度という。$\int_\R f_X=1$ であり、可積分な $g(X)$ に対し
$$\E[g(X)]=\sum_xg(x)p_X(x)\quad\text{または}\quad
  \E[g(X)]=\int_\R g(x)f_X(x)\,dx.$$
密度の高さは一点の確率ではない。密度があれば $P(X=x)=0$ である。
また、分布関数が連続でも密度をもつとは限らない。

\topic{ベルヌーイ分布と二項分布}
$0<p<1$ とし、$P(B=1)=p$、$P(B=0)=1-p$ なら $B\sim\mathrm{Bern}(p)$ と書く。
独立な $B_1,\ldots,B_n\sim\mathrm{Bern}(p)$ の和 $S_n$ は二項分布に従う。
$$
  P(S_n=k)=\binom nkp^k(1-p)^{n-k}\quad(0\le k\le n),
  \qquad \E[S_n]=np,\quad\Var(S_n)=np(1-p).
$$
$k$ 個の成功の配置が $\binom nk$ 通りある。
期待値は線形性から、分散は独立性と $B_j^2=B_j$ から計算できる。

\topic{ポアソン分布とまれな事象の極限}
$\lambda>0$ に対し、$X\sim\mathrm{Poi}(\lambda)$ とは
$$
  P(X=k)=e^{-\lambda}\frac{\lambda^k}{k!}\quad(k=0,1,\ldots),
  \qquad \E[X]=\lambda,\quad\Var(X)=\lambda
$$
である。$\sum_{k\ge0}\lambda^k/k!=e^\lambda$ より確率の総和は $1$。
また $\E[X(X-1)]=\lambda^2$ を使うと分散を求めやすい。
固定した $k$ に対し
$$\binom nk(\lambda/n)^k(1-\lambda/n)^{n-k}
  \longrightarrow e^{-\lambda}\frac{\lambda^k}{k!}.$$
多数の独立な機会に小さな確率で起こる事象の総回数を、この極限が記述する。

\topic{幾何分布：最初の成功までの試行回数}
$X\sim\mathrm{Geo}(p)$ は、本資料では $1,2,\ldots$ の値をとる規約とする。
$$P(X=k)=(1-p)^{k-1}p,\qquad \E[X]=\frac1p,\qquad
  \Var(X)=\frac{1-p}{p^2}.$$
$P(X>m)=(1-p)^m$ より、整数 $m,n\ge0$ に対し
$$P(X>m+n\mid X>m)=\frac{(1-p)^{m+n}}{(1-p)^m}=P(X>n).$$

\newpage
\topic{一様分布：区間の長さに比例する確率}
$a<b$ に対し、$X\sim\mathrm U(a,b)$ の密度とモーメントは
$$
 f_X(x)=\frac{\ind_{[a,b]}(x)}{b-a},\qquad
 \E[X]=\frac{a+b}{2},\qquad \Var(X)=\frac{(b-a)^2}{12}.
$$
$a\le s<t\le b$ なら $P(s<X\le t)=(t-s)/(b-a)$。
$U\sim\mathrm U(0,1)$ なら $a+(b-a)U\sim\mathrm U(a,b)$ である。

\topic{指数分布：連続時間の待ち時間}
$\lambda>0$ を rate とする $X\sim\mathrm{Exp}(\lambda)$ の密度は
$$f_X(x)=\lambda e^{-\lambda x}\ind_{[0,\infty)}(x),\qquad
  P(X>t)=e^{-\lambda t}\quad(t\ge0).$$
部分積分、または尾確率の積分によって
$$\E[X]=\frac1\lambda,\qquad \E[X^2]=\frac2{\lambda^2},\qquad
  \Var(X)=\frac1{\lambda^2}.$$
$s,t\ge0$ に対して $P(X>s+t\mid X>t)=e^{-\lambda s}$。
幾何分布と同じく記憶なし性質をもつが、変数は連続時間を表す。

\topic{ガンマ分布：形状と rate を区別する}
$\alpha,\beta>0$ に対し $\mathrm{Gamma}(\alpha,\beta)$ を rate 表示で定義する。
$$
 \Gamma(\alpha)=\int_0^\infty x^{\alpha-1}e^{-x}\,dx,\qquad
 f_X(x)=\frac{\beta^\alpha}{\Gamma(\alpha)}x^{\alpha-1}e^{-\beta x}
 \quad(x>0).
$$
$x\le0$ では密度を $0$ とする。$\Gamma(\alpha+1)=\alpha\Gamma(\alpha)$ を用いると
$$\E[X^r]=\frac{\Gamma(\alpha+r)}{\beta^r\Gamma(\alpha)}\quad(r> -\alpha),
 \qquad \E[X]=\frac\alpha\beta,\quad\Var(X)=\frac\alpha{\beta^2}.$$
$\alpha=1$ は指数分布である。独立な $\mathrm{Exp}(\beta)$ 変数 $m$ 個の和は
$\mathrm{Gamma}(m,\beta)$ に従い、$m$ 回目の到着までの待ち時間を表す。
scale $\theta=1/\beta$ を使う本では平均は $\alpha\theta$、分散は $\alpha\theta^2$ となる。

\topic{正規分布：中心と幅を二つのパラメータで表す}
$m\in\R$、$\sigma>0$ に対し、$X\sim N(m,\sigma^2)$ の密度は
$$f_X(x)=\frac1{\sqrt{2\pi}\sigma}
 e^{-(x-m)^2/(2\sigma^2)},\qquad \E[X]=m,\quad\Var(X)=\sigma^2.$$
$Z\sim N(0,1)$ なら $m+\sigma Z\sim N(m,\sigma^2)$ である。
対称性から $\E[Z]=0$、部分積分から $\E[Z^2]=1$。
正規化に使う Gauss 積分と、二項分布との関係は次回に扱う。

\newpage
\topic{演習問題}
\begin{enumerate}[label=\arabic*.]
\item $X\sim\mathrm{Exp}(\lambda)$ の平均、二次モーメント、分散を積分で求め、
記憶なし性質を示せ。
\item 分布関数 $F_X$ が連続なら $P(X=x)=0$（すべての $x$）を示せ。
さらに、この性質だけから確率密度の存在を結論してよいか考えよ。
\item $X\sim\mathrm{Gamma}(\alpha,\beta)$ の密度の積分が $1$ であることと、
平均・分散の公式を、変数変換 $u=\beta x$ によって確かめよ。
\item 正の確率変数 $T$ の生存関数 $G(t)=P(T>t)$ が連続で正、
$G(s+t)=G(s)G(t)$（$s,t\ge0$）、$G(0)=1$、$G(t)\to0$ を満たすとする。
$G(t)=e^{-ct}$（ある $c>0$）を示せ。
\end{enumerate}

\topic{レポート問題}
\begin{enumerate}[label=\arabic*.]
\item $X\sim\mathrm{Poi}(\lambda)$ の平均と分散を求めよ。
独立な $X\sim\mathrm{Poi}(\lambda)$、$Y\sim\mathrm{Poi}(\mu)$ に対し、
畳み込み和から $X+Y\sim\mathrm{Poi}(\lambda+\mu)$ を示せ。
\item $X\sim\mathrm{Geo}(p)$ の平均と分散、記憶なし性質を導け。
続いてガンマ分布の応用として、独立な指数待ち時間の和を考え、
何が観測量で、$\alpha$ と $\beta$ が何を表すかを説明せよ。
\end{enumerate}

\topic{考えるための手がかり}
ポアソン分布では $k(k-1)$ を掛けると添字をずらしやすい。
幾何分布では $\sum_{k\ge0}q^k=(1-q)^{-1}$ とその微分、
または $\E[X]=\sum_{m\ge0}P(X>m)$ を使える。
記憶なし性質では、条件付き確率の分母が正であることも確認しよう。
\vfill
{\small\color{muted}解答メモ：パラメータの範囲、値域、密度が $0$ となる範囲を省略しない。}

\columnstart{回数と待ち時間は、同じ過程の二つの見方}
一定時間内に何回起こるかを数えるとポアソン分布、次に起こるまで待つと指数分布。
この関係は分布の公式の偶然の一致ではなく、同じランダムな到着過程を
異なる方向から観測した結果として理解できる。

\columnheading{短い時間での仮定を置く}
$N(t)$ を時刻 $t$ までの到着回数とし、$N(0)=0$ とする。
異なる時間区間の増分が独立で、増分の分布は区間の長さだけで決まると仮定する。
さらに $h\downarrow0$ で
$$P(N(h)=1)=\lambda h+o(h),\qquad P(N(h)\ge2)=o(h)$$
と置く。$\lambda>0$ は単位時間当たりの到着の rate である。
これらの仮定の下で、微小時間に複数回到着する寄与は一次の項より小さい。

\columnheading{回数の確率が満たす方程式}
$p_k(t)=P(N(t)=k)$ とすると、独立増分を使った畳み込みと上の展開から
$$p_0'(t)=-\lambda p_0(t),\qquad
  p_k'(t)=-\lambda p_k(t)+\lambda p_{k-1}(t)\quad(k\ge1).$$
初期条件 $p_0(0)=1$、$p_k(0)=0$（$k\ge1$）を満たす解は
$$p_k(t)=e^{-\lambda t}\frac{(\lambda t)^k}{k!}.$$
したがって、時刻 $t$ の総回数は $\mathrm{Poi}(\lambda t)$ に従う。

\columnheading{最初の到着まで待つ}
$T_1=\inf\{t\ge0:N(t)\ge1\}$ とする。
「$t$ より長く待つ」は「時刻 $t$ までに一度も到着しない」と同じだから
$$P(T_1>t)=P(N(t)=0)=e^{-\lambda t}.$$
よって $T_1\sim\mathrm{Exp}(\lambda)$ であり、記憶なし性質も直ちに従う。
すでに $t$ 待ったという情報を得ても、さらに待つ時間の分布は同じである。

\columnheading{$m$ 回目の到着まで待つ}
$T_m$ を $m$ 回目の到着時刻とすると
$$P(T_m>t)=P(N(t)\le m-1)
 =e^{-\lambda t}\sum_{k=0}^{m-1}\frac{(\lambda t)^k}{k!}.$$
この式を微分すると $T_m$ の密度は
$\lambda^m t^{m-1}e^{-\lambda t}/(m-1)!$（$t>0$）となる。
これは $\mathrm{Gamma}(m,\lambda)$ である。
どの現象にも当てはまるわけではなく、独立な増分と一定 rate という仮定が
どこまで適切かを確かめることが、分布名を選ぶ前の大切な作業となる。

\label{lastpage}
\end{document}
